% ===================================================================
%  TABELLA nr. 15 — Il martgà. --- Las victualias.
%  Traduction 2026 d'apres texto/io, controlee sur texto/fr, texto/ca
%  et texto/oc. Consigne : voir texto/rm/10-tabella-01.tex.
%
%  LA POMME DE TERRE DE L'ALINEA 10 CLOT UNE SERIE COMMENCEE DEUX
%  COLONNES PLUS TOT, ET C'EST ELLE QUI TRANCHE.
%
%  Le tableau 15 lituanien avait conclu de cet etal qu'un legume
%  arrive apres la fixation du vocabulaire se reconnait a ce qu'il
%  est emprunte. Le tableau 15 luxembourgeois avait corrige : la date
%  fixe qu'il y aura un mot neuf, elle ne fixe pas lequel. Six
%  langues, six reponses :
%
%    lituanien      « bulvė »     — un emprunt ;
%    tcheque        « brambory »  — le nom du pays d'ou elle vient ;
%    allemand       « Kartoffel » — le nom italien de la truffe ;
%    francais       « pomme de terre » — un compose fait chez soi ;
%    luxembourgeois « Gromper »   — un autre compose fait chez soi ;
%    ROMANCHE       « tartuffel » — LE MOT ALLEMAND.
%
%  ET C'EST LA QUE LA SERIE SE DECIDE. Le romanche est une langue
%  ROMANE ; son voisin roman, l'italien, dit « patata ». Il aurait
%  suffi de tendre la main au sud. Le romanche a pris le mot du nord,
%  parce que la pomme de terre lui est arrivee du nord, par les cols,
%  avec les semences et le mode de culture.
%
%  C'est exactement la regle du tableau 10, verifiee cette fois sur
%  une plante et non sur un gouvernail : QUAND ON EMPRUNTE, C'EST
%  L'ORIGINE DE LA CHOSE QUI CHOISIT LE PRETEUR, PAS LA PARENTE DES
%  LANGUES. Le luxembourgeois disait que la date ne decide pas de la
%  forme du mot ; le romanche ajoute que ce qui decide, quand la
%  forme est un emprunt, c'est la route de la chose.
%
%  On note pour etre juste que plusieurs idiomes des vallees disent
%  aussi « pom da terra », compose sur le modele francais. Le fonds
%  commun, lui, est bien la forme du nord.
% ===================================================================

%%K t15-apar-1 apar
\VUsaut{4.0mm}
\VUtitre{15.0pt}{60}{TABELLA nr. 15}
\VUsaut{3.0mm}

%%K t15-tit-1 sub
\VUpk{11.4pt}{40}{Il martgà. --- Las victualias.}
\VUsaut{3.0mm}

%%K t15-01-1 p
1. --- Ils pli blers commerziants che furneschan a nus la rauba necessaria per noss
nutriment sa radunan sut la \VUgras{halla} \textsuperscript{(1)} dal \VUgras{martgà cuvert} u en
las vias vischinantas.

%%K t15-02-1 p
2. --- Il \VUgras{mazlader da portgs} \textsuperscript{(2)}, che venda \VUgras{schambun}
\textsuperscript{(3)}, \VUgras{salsiz da sang} \textsuperscript{(4)}, \VUgras{salsiz}
\textsuperscript{(5)}, \VUgras{salsiz da tripa} \textsuperscript{(6)} e \VUgras{lardun}
\textsuperscript{(7)}, stat sper la \VUgras{vendidra da paintg e chaschiel}
\textsuperscript{(8)}, da la quala l'étalage è pervis cun \VUgras{chaschiels ollandais}
\textsuperscript{(9)} cun crusta cotschna, cun \VUgras{chaschiels Camembert}
\textsuperscript{(10)}, cun \VUgras{rodas da Gruyère} \textsuperscript{(11)} e cun frais
\VUgras{tocs da paintg} \textsuperscript{(12)}.

%%K t15-03-1 p
3. --- Avant ella porscha la \VUgras{vendidra da verduras} \textsuperscript{(13)} a sias
clientas bellas \VUgras{tomatas} \textsuperscript{(14)}, \VUgras{cabis-flur}
\textsuperscript{(15)}, \VUgras{salata} \textsuperscript{(16)}, \VUgras{chaus da cabis}
\textsuperscript{(17)}, \VUgras{cucurbitas} \textsuperscript{(19)}, \VUgras{melonas}
\textsuperscript{(18)} ed era \VUgras{bulieus} \textsuperscript{(20)}. Ma in giuven da la biera,
ch'ha fermà ses \VUgras{char da liuraziun} \textsuperscript{(21)} chargià cun \VUgras{barigls da
biera} \textsuperscript{(22)} gist avant ses étalage, impedescha sias clientas da
s'avischinar. Ina giardiniera al porta in chanaster d'\VUgras{artischoccas}
\textsuperscript{(23)}.

%%K t15-tit-2 sub
\VUsaut{4.0mm}
\VUpk{10.2pt}{40}{Vender e cumprar.}
\VUsaut{3.0mm}

%%K t15-04-1 p
4. --- Sin il martgà è l'animaziun gronda, perquai ch'l'ura da la \VUgras{vendita a
l'incant} \textsuperscript{(24)} è arrivada. Las \VUgras{patrunas da chasa}
\textsuperscript{(26)}, che vegnan là per lur cumpras, restan en passond avant la buttega
da la \VUgras{vendidra da tripa} \textsuperscript{(25)}, per cumprar \VUgras{tripas} u in
\VUgras{chau da vadè}.

%%K t15-05-1 p
5. --- In grass \VUgras{cuschinunz} \textsuperscript{(27)} martgadegia in fitg bel
\VUgras{tonn} tar la \VUgras{vendidra da pesch} \textsuperscript{(28)}, che auza ses pretschs a
bel plaschair. Ella ha survegnì ina gronda quantitad d'\VUgras{anguillas, solas,
truttas} e \VUgras{carpas}. Ses \VUgras{merluz} e sias \VUgras{mislas} èn fitg frais.

%%K t15-tit-3 sub
\VUsaut{4.0mm}
\VUpk{10.2pt}{40}{Rauba bunmartgada e chara.}
\VUsaut{3.0mm}

%%K t15-06-1 p
6. --- La \VUgras{vendidra da pulaglia} \textsuperscript{(29)}, installada sper el, è fitg
conscienziusa. Cur ch'ella venda in \VUgras{tatschin}, ina \VUgras{aucha}, ina
\VUgras{anda} u in simpel \VUgras{pulschin}, pretenda ella mo il pretsch endretg.

%%K t15-07-1 p
7. --- A l'chantun da la plazza, sin l'emprim plaun, è dapi pauc temp installà in
\VUgras{vendider da telas} \textsuperscript{(30)}, tar il qual las cumpradras èn adina segiras
da chattar in fitg bel assortiment da \VUgras{tavatschas}, \VUgras{serviettas},
\VUgras{scussals} e \VUgras{sfruschamauns}. Sin il medem plaun ha in \VUgras{curtellier}
\textsuperscript{(31)} installà ses atelier. El agiavgia ils \VUgras{curtels} da las
vendidras en la halla e fa per ils giardiniers \VUgras{seppas} da l'\VUgras{atschal} il pli
dir.

%%K t15-tit-4 sub
\VUsaut{4.0mm}
\VUpk{10.2pt}{40}{La spezeria.}
\VUsaut{3.0mm}

%%K t15-08-1 p
8. --- Sut ses atelier è dapi pauc temp avert in grond magasin da victualias. El
n'è gia betg pli la simpla e nunimpurtanta \VUgras{spezeria} \textsuperscript{(32)} d'ina giada,
che vendeva mo la rauba dal diever cuminaivel, \VUgras{te} \textsuperscript{(33)},
\VUgras{tschigulatta} \textsuperscript{(34)}, \VUgras{zucker en toc} \textsuperscript{(37)} e
\VUgras{savun} \textsuperscript{(38)}. Là venda ins mintga rauba, \VUgras{biscuits crocants},
\VUgras{conservas} \textsuperscript{(35)}, \VUgras{liquors} \textsuperscript{(36)}, \VUgras{rum},
\VUgras{cognac}, \VUgras{kirsch}, \VUgras{anis} e \VUgras{curaçao}. Là chatta ins
selvaschina: \VUgras{levra} \textsuperscript{(39)}, \VUgras{tordas} \textsuperscript{(40)},
\VUgras{pernischs} \textsuperscript{(41)}, \VUgras{quaglias} \textsuperscript{(42)}, \VUgras{anda
selvadia} \textsuperscript{(43)} u \VUgras{fasan} \textsuperscript{(44)}, precis uschè
facilmain sco
fritgas exoticas, sco \VUgras{bananas} \textsuperscript{(46)} ed \VUgras{ananas}.

%%K t15-09-1 p
9. --- In \VUgras{giuven} \textsuperscript{(45)} dal spezier, fitg serviziaivel, è pront a dar
quai ch'ins giavischa: \VUgras{confitira} \textsuperscript{(47)} da ribas u da cotogns,
\VUgras{grascha} \textsuperscript{(48)}, \VUgras{cornichons} \textsuperscript{(49)} u
\VUgras{sardinas}
\textsuperscript{(50)} salidas.

%%K t15-09-2 p
Quai è fitg cumadaivel per las \VUgras{clientas} \textsuperscript{(51)}, che chattan, sper la
rauba la pli communa, las primurs las pli raras e las fritgas las pli savurusas:
succus \VUgras{pirs} \textsuperscript{(52)}, \VUgras{brognaclas} \textsuperscript{(53)},
\VUgras{rapas}
\textsuperscript{(54)}, \VUgras{persis} \textsuperscript{(55)}, \VUgras{mandels}
\textsuperscript{(56)} e \VUgras{nuschs} \textsuperscript{(57)}.

%%K t15-tit-5 sub
\VUsaut{4.0mm}
\VUpk{10.2pt}{40}{La clientella.}
\VUsaut{3.0mm}

%%K t15-10-1 p
10. --- Il \VUgras{ris} \textsuperscript{(58)} e las \VUgras{lantschas} \textsuperscript{(59)}
èn sper
la chascha dals \VUgras{arengs} \textsuperscript{(60)} fimentads; ils \VUgras{tartuffels}
\textsuperscript{(61)} e las \VUgras{fascholas} \textsuperscript{(63)} èn vischins als
\VUgras{mails}
\textsuperscript{(62)}, a las \VUgras{feias} \textsuperscript{(64)}, a las \VUgras{brognaclas
setgas} \textsuperscript{(65)} ed a las \VUgras{nuschellas} \textsuperscript{(66)}. Ins chatta en
quel magasin frais \VUgras{ovs} \textsuperscript{(67)} ed \VUgras{olivas} \textsuperscript{(68)};
perquai vegnan \VUgras{chanasters} \textsuperscript{(70)} e \VUgras{raits da provisiun}
\textsuperscript{(73)} emplenids fitg svelt.

%%K t15-11-1 p
11. --- Il \VUgras{caffè} \textsuperscript{(72)} da quella excellenta firma ha survegnì,
tranter las \VUgras{serventas} \textsuperscript{(69)} e las \VUgras{cuschinunzas}, ina fama
merteivla, perquai ch'il vegl \VUgras{spezier} \textsuperscript{(71)} sez al brascha cun la pli
gronda premura.

%%K t15-tit-6 sub
\VUsaut{4.0mm}
\VUpk{10.2pt}{40}{Ils spectaculs.}
\VUsaut{3.0mm}

%%K t15-12-1 p
12. --- Betg tuts van al martgà per sa pervair. En la pittoresca circulaziun sin la
vietta che maina a la halla vesa ins las persunas las pli differentas. Sper in
amiaivel \VUgras{servent dal mazlader} \textsuperscript{(75)}, che tira avant sai ina
\VUgras{carriola} \textsuperscript{(74)}, qua èn dus schuldads en congedi, fitg fiers da
mussar lur uniforma splendurusa: il \VUgras{hussar} \textsuperscript{(76)} cun blau
\VUgras{dolman} e \VUgras{chako cun plimas}, ed il \VUgras{caporal dals zuavs}, cun
chautschas sfunsadas ed il chau cuvrì cun in \VUgras{fez}.

%%K t15-13-1 p
13. --- Il \VUgras{pastrizer} \textsuperscript{(80)}, ch'è en pes sin il sulier da la
\VUgras{pastizaria} \textsuperscript{(78)}, ha gist impastà la pasta da ses \VUgras{paun}
\textsuperscript{(79)} ed ha prendì ord il fuorn ils \VUgras{cornets} \textsuperscript{(81)}, ils
\VUgras{panets} \textsuperscript{(82)} ed ils \VUgras{pauns tunds} \textsuperscript{(83)} ch'èn
en la
vitrina.

%%K t15-14-1 p
14. --- Sur la pastizaria abitescha ina abila \VUgras{biancariera} \textsuperscript{(84)}, che
dat lavur a pliras \VUgras{bruderaras} \textsuperscript{(85)}. Sper ella abitescha ina
\VUgras{lavandera-stirunza} \textsuperscript{(86)}, che amidunescha la \VUgras{biancaria}
\textsuperscript{(88)} suenter l'avair lavada e sitgentada, e che la stira cun in
\VUgras{fier da stirar} \textsuperscript{(87)}.

%%K t15-tit-7 sub
\VUsaut{4.0mm}
\VUpk{10.2pt}{40}{Charn, peschs e verduras.}
\VUsaut{3.0mm}

%%K t15-15-1 p
15. --- En la \VUgras{mazlaria} \textsuperscript{(89)}, situada sin il plaunterren, venda ins
mintga sort da \VUgras{charn}, \VUgras{charn da muntun} \textsuperscript{(90)}, \VUgras{charn da
bov} \textsuperscript{(94)} u \VUgras{charn da vadè} \textsuperscript{(92)}, sco \VUgras{cossas da
muntun, beefsteaks, rustids} e \VUgras{cotletas}. Il \VUgras{giuven dal mazlader}
\textsuperscript{(93)} disseca tut quellas charns e metta da la vart las \VUgras{ossa cun
mieul} \textsuperscript{(96)}. Vi da la porta da la mazlaria è ina \VUgras{vendidra d'ustras}
\textsuperscript{(97)}, che venda \VUgras{ustras} \textsuperscript{(98)}. Ella las avra sch'ins
al giavischa.

%%K t15-16-1 p
16. --- Tut quels commerziants pajan patenta u taxa da plaz; perquai persequitescha
il \VUgras{polizist} \textsuperscript{(100)} senza cumpassiun las paupras \VUgras{colportadras da
verduras} \textsuperscript{(99)}, che als fan concurrenza, purtond enturn cun lur charrettas
\VUgras{carottas, ravanels, raps, porrs} u \VUgras{faschs da sparsla, pisels frais,
spinat, acetusa, crescha} e \VUgras{petersiglia}.

%%K t15-tit-8 sub
\VUsaut{4.0mm}
\VUpk{10.2pt}{40}{Fai attenziun al polizist!}
\VUsaut{3.0mm}

%%K t15-17-1 p
17. --- Il mat che venda \VUgras{agl} \textsuperscript{(101)} e \VUgras{tschagulas} sa fitg bain
ch'era el na dastga betg vender sia rauba sper il martgà. El fugia, currind cun il
ristg da rebaltar il chanaster da \VUgras{granzs}, \VUgras{granzs d'aua dultscha} e
\VUgras{crevettas} dal \VUgras{vendider} \textsuperscript{(102)} da \VUgras{langustas} e
\VUgras{homards}.

%%K t15-18-1 p
18. --- Tuts sa movan e fan canaster; ma quai n'impedescha betg d'udir cleramain la
vusch dal \VUgras{vaidrer} \textsuperscript{(103)} ambulant, ni il clom dal \VUgras{carbunier}
\textsuperscript{(104)}, ch'ha gist bandunà per intgins mumaints sia charretta per liurar
in \VUgras{satg da carvun} \textsuperscript{(105)}.

\VUsaut{6.0mm}
\VUfilet{22mm}
